% =====================================================================
%  APÊNDICE B — GUIA DE INSTALAÇÃO E EXECUÇÃO
% =====================================================================

\chapter{GUIA DE INSTALAÇÃO E EXECUÇÃO}
\label{ap:instalacao}

Este apêndice apresenta os pré-requisitos e os passos para instalação e
execução do \argus{} em modo de desenvolvimento.

\section{Pré-requisitos}
\label{ap:instalacao-pre-requisitos}

\begin{itemize}
    \item Go (v1.20 ou superior): \url{https://go.dev/};
    \item Node.js e npm;
    \item Wails CLI: \code{go install github.com/wailsapp/wails/v2/cmd/wails@latest};
    \item Android Studio (apenas para a compilação móvel via Capacitor).
\end{itemize}

\section{Instalação das dependências}
\label{ap:instalacao-dependencias}

\begin{enumerate}
    \item Clonar o repositório em
        \url{https://github.com/SoNdA11/argus-cyclist} e entrar na pasta do
        projeto;
    \item Instalar as dependências do \textit{backend} Go:
        \begin{verbatim}
        go mod tidy
        \end{verbatim}
    \item Instalar as dependências do \textit{frontend}:
        \begin{verbatim}
        cd frontend
        npm install
        cd ..
        \end{verbatim}
\end{enumerate}

\section{Execução em modo de desenvolvimento}
\label{ap:instalacao-execucao}

Com as dependências instaladas, a aplicação pode ser executada em modo de
desenvolvimento (com \textit{hot reload}) por meio do comando:

\begin{verbatim}
wails dev
\end{verbatim}

\section{Geração do aplicativo de desktop}
\label{ap:instalacao-build}

Para gerar o binário de produção para Windows ou Linux, execute, na raiz do
projeto:

\begin{verbatim}
wails build
\end{verbatim}

O artefato é criado no diretório \code{build/bin}. A compilação para macOS
depende de macOS com Xcode e não foi validada neste trabalho.

\section{Compilação para Android}
\label{ap:instalacao-android}

A versão Android reutiliza o \textit{frontend} e o empacota com o Capacitor.
Após instalar o Android Studio e configurar o SDK, execute:

\begin{verbatim}
cd frontend
npm run build
npx cap sync android
npx cap open android
\end{verbatim}

O último comando abre o projeto no Android Studio, a partir do qual o APK pode
ser gerado e instalado em um dispositivo físico. O projeto iOS/iPadOS existe,
mas não foi compilado nem validado neste trabalho.
